\clearpage
\section{Conclusions}

This project implements the three required execution strategies without
turning them into three separate codebases. The concurrent versions share one
\texttt{GameLoop}, one command protocol, and one
\texttt{ParallelPhysicsEngine}. Their only relevant difference is the
\texttt{RangeScheduler}: persistent platform threads in one case, executor
tasks in the other.

The controller remains the single writer of the game protocol. Physics workers
operate only inside bounded phases, on disjoint ranges or private
accumulators, and the deterministic merge completes before publication. This
is the central trade-off of the solution: parallelize the expensive local
work, while keeping the global meaning of a tick simple and explicit.

The benchmark evidence in the repository and the rendered figures show consistent but limited gains. For the
speedup workloads, both concurrent versions beat the sequential baseline from
1000 balls upward, and the best median in that set is the executor run on 2500
balls (\texttt{242.8611 ms}, \texttt{1.819991} speedup). For the 2500-ball
scalability run, both implementations reach their best measured point at
\textbf{16 workers}, and performance then flattens or slightly regresses at 17,
so scaling remains \textbf{sublinear}.

The verification evidence is intentionally narrower than the full game. The
Petri nets and JPF harnesses model the synchronization protocol, not numerical
physics or performance. Together with the JUnit suite and reproducible
benchmarks, they provide evidence at the right level for each claim: functional
correctness, absence of errors in the bounded protocol, and measured---rather
than assumed---parallel speedup.
